\chapter{Noncircularity and Dependency Audit}
\label{app:perverse-dependency-audit}

This appendix makes the proof graph explicit.  An arrow means that the target
uses the source.  There is no implicit appeal to a theorem later in the graph.

\section{The main graph}

\[
\begin{tikzcd}[column sep=small,row sep=large]
\substack{\text{six operations}\\\text{base change, trace, duality}}
 \arrow[r] \arrow[dr] &
\substack{\text{perverse \(t\)-structure}\\j_{!*},\ \IC,\ \text{Artin vanishing}}
 \arrow[d] \\
&
\substack{\text{Fourier inversion}\\\text{\(t\)-exactness, Plancherel}}
 \arrow[d]\\
\substack{\text{rank-one purity}\\\text{real-sheaf theorem}}
 \arrow[r]&
\substack{\text{curve weight estimate}\\H_c^a\text{ upper weights}}
 \arrow[d]\\
\substack{\text{Noether normalization}\\\text{excision, Leray filtration}}
 \arrow[r]&
\substack{Rf_!\text{ preserves}\\\text{upper weights}}
 \arrow[d]\\
\substack{\text{Verdier duality}\\Rf_!=Rf_*\text{ for proper }f}
 \arrow[r]&
\substack{\text{proper purity}\\\text{IC/IH purity}}
 \arrow[d]\\
\substack{\text{trace and Poincare}\\\text{determinant algebra}}
 \arrow[r]&
\substack{\text{smooth-proper purity}\\\text{zeta and functional equation}}
 \arrow[d]\\
\substack{\text{universal-\(\iota\) weights}\\\text{Fatou and Gauss}}
 \arrow[r]&
\substack{P_i\in\mathbf Z[T]\\\ell\text{-independence}.}
\end{tikzcd}
\tag{B.1}\label{eq:dependency-graph}
\]

The decomposition theorem branches from proper purity and points forward:
\[
  \text{proper purity}
  \Longrightarrow
  \text{pure perverse semisimplicity}
  \Longrightarrow
  \text{pure-complex splitting}
  \Longrightarrow
  \text{decomposition}.
\tag{B.2}\label{eq:decomposition-branch}
\]
There is no arrow from \eqref{eq:decomposition-branch} back to the curve or
direct-image theorem.

\section{Audit of tempting circular shortcuts}

\begin{longtable}{
 >{\raggedright\arraybackslash}p{0.31\textwidth}
 >{\raggedright\arraybackslash}p{0.30\textwidth}
 >{\raggedright\arraybackslash}p{0.29\textwidth}}
\toprule
Tempting shortcut & Why it is circular or false & Replacement used here\\
\midrule
\endhead
“\(Rf_*\) of a pure complex is pure by BBD.” &
BBD proper purity uses Deligne's direct-image weight theorem, whose
smooth-proper specialization is the RH being proved. &
Prove the curve estimate, globalize \(Rf_!\) upper weights, then dualize.\\

“The decomposition theorem proves smooth-proper purity.” &
The finite-field decomposition theorem is proved from purity and weight
separation. &
Purity is proved before decomposition; the latter is recorded as a
corollary.\\

“Fourier transform raises weights by one.” &
The Fourier transform is an \(Rp_!\) construction, so this assertion is a
case of the direct-image theorem. &
Before the curve theorem use only stalk formulas, inversion,
\(t\)-exactness, real mixedness, and Plancherel.\\

“The perverse spectral sequence degenerates.” &
No such general degeneration is available at that stage. &
Use only its finite Frobenius-stable filtration and subquotients.\\

“Rationality of \(Z\) makes each \(P_i\) rational.” &
Alternating numerator and denominator factors can cancel. &
First prove universal weights; distinct moduli rule out cancellation.\\

“Pure perverse means Frobenius semisimple.” &
Purity controls eigenvalue moduli, not Jordan blocks. &
State only geometric semisimplicity after forgetting Frobenius.\\

“\(\Qbarl[d]\) is perverse on every pure \(d\)-fold.” &
Singularities may violate cosupport. &
Use \(\IC_X\); identify it with \(\Qbarl[d]\) only when \(X\) is smooth.\\
\bottomrule
\end{longtable}

\section{The weight-shift audit}

All shifts can be reconstructed from four tests:
\[
\begin{array}{c|c}
\text{object}&\text{weight}\\
\hline
\Qbarl&0\\
\Qbarl(1)&-2\\
\cL[d]\text{ for \(\cL\) pure \(u\)}&u+d\\
\pH^a(\cK)\text{ for \(\cK\) pure \(w\)}&w+a.
\end{array}
\tag{B.3}\label{eq:weight-shift-audit}
\]
For \(X\) smooth proper of dimension \(d\),
\[
  \Qbarl[d]\xrightarrow{Ra_*}
  Ra_*\Qbarl[d]
\xrightarrow{H^{\,i-d}}
  H^i(X,\Qbarl)
\]
has successive weight labels
\[
  d,\qquad d,\qquad d+(i-d)=i.
\tag{B.4}\label{eq:terminal-shift-audit}
\]
Any proof producing \(i-d\), \(i+d\), or \(2d-i\) at the last stage has lost a
shift or dualized the wrong object.

\section{Frobenius audit}

The notes use geometric Galois Frobenius on stalks and cohomology.  Therefore
\[
 F|\Qbarl(1)=q^{-1},\qquad F|\Qbarl(-1)=q,
\]
\[
 P_i(T)=\det(1-TF\mid H^i),\qquad
 |\iota(\alpha)|=q^{i/2}
\]
for an eigenvalue \(\alpha\) of \(F\).  The polynomial roots in the variable
\(T\) are \(\alpha^{-1}\).  Poincare duality pairs eigenvalues
\(\alpha\) and \(q^d/\alpha\).  These conventions are tested on
\(\PP^n\), whose zeta function is
\[
  Z(\PP^n,T)=\prod_{r=0}^{n}(1-q^rT)^{-1}.
\]

\section{Use of all complex realizations}

For each chosen
\(\iota:\Qbarl\simeq\mathbf C\), the entire weight proof yields
\(\iota\)-purity.  Only after repeating it for all \(\iota\) do we use the
arithmetic argument
\[
\begin{gathered}
Z\in\mathbf Z[[T]]\cap\mathbf Q(T)\\
\Downarrow\\
\text{integral reduced numerator and denominator}\\
\Downarrow\quad\text{universal modulus separates weights}\\
P_i(T)\in\mathbf Z[T]\text{ and is independent of }\ell.
\end{gathered}
\]
A fixed embedding gives an analytic modulus but not Galois stability of the
weight strata.

\section{Terminal theorem audit}

The final theorem depends on the following proved objects, not theorem
constants introduced as assumptions:

\begin{enumerate}
\item finite-dimensional \(H^i\) with geometric Frobenius;
\item trace formula for all powers of Frobenius;
\item Poincare duality with multiplier \(q^d\);
\item the derived \(Rf_!\) upper-weight theorem;
\item its Verdier-dual lower theorem and proper-purity corollary;
\item the perverse identity \(\IC_X=\Qbarl[d]\) for smooth \(X\);
\item the algebraic integral-factor and separation lemmas.
\end{enumerate}

Deleting the chapters on semisimplicity and decomposition would not break the
classical Weil theorem.  Deleting proper purity would.  This is the simplest
mechanical check that decomposition has not been used circularly.
